\FloatBarrier
\section{本章小结}
\label{sec:solution-error-chapter-summary}

本章把一个求解结果的可信程度拆成三种来源分别处理。第一种是问题本身的敏感性：可逆方程的解在右端或系数矩阵受有限扰动时，相对误差不超过条件数乘相对扰动，再除以到奇异边界的余量$1-\kappa(\bm A)\rho_A$；误差沿小奇异值方向放大最多。

最小二乘中，拟合值与残差的变化受观测变化控制，系数却按$1/\sigma_r$放大，拟合得好不等于系数可靠。伪逆只在正奇异值有共同下界时连续，秩增加处不连续，所以数值秩要按误差水平而不是按是否精确为零来确定。

第二种是算法在有限精度下引入的误差。残差小只说明后向误差小；前向误差约不超过条件数乘后向误差。后向稳定的算法也不能改善病态问题，但能避免额外损失：形成正规方程使条件数平方，而QR、SVD与LSQR、LSMR直接作用于$\bm A$。第三种是迭代提前停止留下的误差，需要用真实残差和相应的正规残差核验，并分清达到容差与只是停止计算。

当问题本身太敏感时，可以主动改变目标。Tikhonov正则化等价于增广最小二乘，解唯一当且仅当$\ker\bm A\cap\ker\bm\Gamma=\{\bm0\}$；岭回归在奇异方向上把增益$1/\sigma_i$换成有界的$\sigma_i/(\sigma_i^2+\lambda)$。谱截断只在保留的奇异方向上取倒数，是限制子空间中的最小二乘解，误差分成保留方向的噪声放大与舍弃方向的偏差两部分。两者都以拟合误差和偏差换取稳定，参数要联系噪声规模与任务确定；它们解的是新目标，较小的系数不能当作原问题解得更准。

选择方法时，先固定求解目标与验收标准：相容方程看原残差，不相容最小二乘还要看正规残差，最小范数还要排除零空间分量。再按矩阵形状、对称正定性、秩、稀疏性与数据访问方式确定可用的方法，以相同精度下的准备、单步、步数、存储与通信总成本比较，并分开问题病态、算法误差与提前停止三种误差。

没有一种方法在所有问题上都更好：直接分解适合能存储因子或多右端的问题，Krylov方法适合大型稀疏问题，LSQR、LSMR保留矩形最小二乘的结构，行投影法适合逐行访问的数据。
